% ASSERT_CONVENTION: natural_units=natural, metric_signature=mostly_minus,
%   jordan_product=(1/2)(ab+ba), octonion_basis=fano_e1e2=e4,
%   complex_structure=u_equals_e7, primitive_idempotent=E11_diag_1_0_0,
%   det_ssot=ring_lemma_det_3, qgt=Tr(P_dP_dP)_over_Q(i)_i_symbolic,
%   berry_curvature_FB=minus2_Im_QGT, fubini_study=Re_QGT, cosmological_constant=Lambda_0
%
% Phase 76, Plan 01: Phase A.5 -- THE BERRY-CURVATURE SAME-WALL GATE (SOFT KILL, v18.0)
%   PART 1 (the VACUUM half).  Conjunctive with Phase 75 (SURVIVES).
%
% THIS PLAN (76-01) WRITES PART 1 ONLY:
%   (a) the QGT recipe + the CP^1 convention pin (F_B = -sin theta/2, g_thth = 1/4);
%   (b) VALD-03: F_B = Im(QGT) is WELL-DEFINED + generically NONZERO, BORN FROM the
%       C_u breaking (OP^2 = F_4/Spin(9) isotropy-irreducible => no invariant 2-form
%       => round Berry = 0);
%   (c) CALC-03 (SOFT, INFORMATIVE, NEVER a KILL): Re(QGT) = a positive-definite
%       Fubini-Study metric; its non-Einstein CHARACTER vs the v17.0 cone-Hessian;
%   (d) CALC-04: the M=0 vacuum class of F_B (pure-Lambda/Kahler).
% The DECISIVE matter-on VALD-04 (Einstein-vs-EM same-wall) verdict + the explicit
% SOFT KILL / SURVIVES line are appended by plan 76-02 (depends_on this plan).
% NO VALD-04 verdict is pre-empted here.
%
% Backed exact-over-Q(i) by code/cartan_phaseA5_berry.py (SymPy 1.14.0; ALL_PASS,
% exit 0).  The imaginary unit i is sympy.I, SYMBOLIC (e_7 -> i); every decisive
% number is rational/symbolic from sympy diff / re / im / .eigenvals over QQ or
% Q(i) -- NEVER numpy/float (fp-float-decisive rejected).  octonion_algebra.py is
% ABSENT on the decisive path (fp-octonion-algebra; runtime source guard).
%
% References:
%   [PV80]  J.P. Provost, G. Vallee, "Riemannian structure on manifolds of quantum
%             states", Comm. Math. Phys. 76 (1980) 289. The QGT split: Re = Fubini-
%             Study metric, Im = Berry curvature. THE load-bearing convention -- the
%             dead route (v17.0) was Re, this route is Im.
%   [Baez]  J.C. Baez, "The Octonions", math/0105155 (2002), Sec. 3.4. OP^2 =
%             F_4/Spin(9); isotropy = the unique irreducible 16-dim Spin(9) spinor
%             rep => isotropy-irreducible => unique invariant symmetric form (FS) and
%             NO invariant 2-form => round Berry = 0. THE born-from-breaking fact.
%   [Berger] M. Berger, classification of isotropy-irreducible homogeneous spaces
%             (an isotropy-irreducible space has a unique invariant metric up to scale).
%   [FK94]  J. Faraut, A. Koranyi, "Analysis on Symmetric Cones" (1994). The cone
%             metric g_X = Hess(-log det) underlying the cone-Hessian diag(9,9,18,18)
%             that CALC-03 SOFT-compares Re(QGT) to (REAL-part character check ONLY).
%   [v17]   code/bulk_geometry_verification.py (cone_hessian_at_center,
%             h3_constant_curvature, ricci_decomposition_n4); v17.0 cone-Hessian NONE.
%             Warm engine for the CALC-03 Re-anchor + the CALC-04 decomposition. The
%             v17.0 NONE binds ONLY the Re sector, NOT this Im verdict.
%   [P46/EMB] Phase 46 / code/embedding_under_E_verification.py: slice_to_complex
%             (e_7 -> i), E, proj_u_exact -- the C_u reduction; the e_7 -> i bridge
%             that makes Im the literal SymPy imaginary part over Q(i).
%   [P75]   derivations/75-coframe-reduction.tex / 75-01-SUMMARY.md -- Phase A
%             SURVIVES: (E_11,u=e_7) forces a 4d Lorentzian (1,3) coframe carrying
%             SO(3,1); the C_u^2 survivors {11,18,19,26} = the 4 base directions here.

\documentclass[11pt]{article}
\usepackage{amsmath,amssymb,amsthm}
\usepackage[margin=1in]{geometry}
\usepackage{hyperref}

\newcommand{\Tr}{\operatorname{Tr}}
\renewcommand{\Re}{\operatorname{Re}}
\renewcommand{\Im}{\operatorname{Im}}
\newcommand{\dd}{\mathrm{d}}
\newcommand{\Cu}{\mathbb{C}_u}
\newcommand{\HO}{\mathfrak{h}_3(\mathbb{O})}
\newcommand{\HCu}{\mathfrak{h}_3(\Cu)}
\newcommand{\OP}{\mathbb{OP}^2}
\newcommand{\Spin}{\mathrm{Spin}}
\newcommand{\SO}{\mathrm{SO}}
\newcommand{\FB}{F_{\mathrm{B}}}

\theoremstyle{plain}
\newtheorem{proposition}{Proposition}
\newtheorem{lemma}{Lemma}
\theoremstyle{definition}
\newtheorem{definition}{Definition}
\newtheorem{remark}{Remark}

\title{Phase A.5 --- The Berry-Curvature Same-Wall Gate (SOFT KILL)\\
\large PART 1: The QGT recipe, well-definedness, and the $M=0$ vacuum class}
\author{v18.0 ``Gravity as the Curvature of the Peirce-Frame Connection on $\HO$''}
\date{Phase 76, Plan 01 (the vacuum half; conjunctive with Phase 75)}

\begin{document}
\maketitle

\begin{abstract}
This plan establishes, exactly over $\mathbb{Q}(i)$ and \emph{before any matter is
turned on}, the vacuum baseline of the canonical Berry curvature
$\FB = \Im(Q) = -2\,\Im\Tr(P\,\dd P\,\dd P)$ of the rank-$1$ $\Cu$-idempotent family
$E(x)$ on the $4$-dimensional slice $\{11,18,19,26\}=\Cu^2$ identified in Phase~75.
We (i) pin the QGT recipe against the $\mathbb{CP}^1$ monopole anchor
($\FB=-\sin\theta/2$, $g_{\theta\theta}=1/4$, flux $=-2\pi$); (ii) prove (VALD-03)
that $\FB$ is well-defined and generically nonzero, \emph{born from the $\Cu$
breaking} of the isotropy-irreducible $\OP=F_4/\Spin(9)$ (which carries no invariant
$2$-form, so the round Berry curvature vanishes identically); (iii) run the soft,
informative real-part anchor (CALC-03) showing $\Re(Q)$ is a positive-definite
Fubini--Study metric, and compare its non-Einstein \emph{character} to the v17.0
cone-Hessian (a discrepancy is expected and is never a kill); and (iv) classify the
$M=0$ vacuum (CALC-04) as pure-$\Lambda$/K\"ahler ($\FB=-2\,\omega_K$, primitive
remainder zero). The decisive matter-on Einstein-vs-EM verdict (VALD-04) and the
explicit \textsc{soft kill}/\textsc{survives} line are deferred to plan 76-02.
\end{abstract}

\section{Setup and the two design points}
\label{sec:setup}

Phase~75 (the coframe-reduction kill gate) \emph{survives}: with the primitive
idempotent $E_{11}=\mathrm{diag}(1,0,0)$ and the complex structure $u=e_7$ fixed, the
$\Cu$ bottleneck $\pi_u$ reduces the $16$-dimensional Peirce half-eigenspace
$V_{1/2}(E_{11})$ to a $4$-dimensional Lorentzian $(1,3)$ coframe carrying
$\SO(3,1)$, forced by $(E_{11},u)$. The four surviving engine directions are
$\{11,18,19,26\}=\Cu^2$. These four directions are the differentiation parameters
$(a,b,c,d)$ throughout this plan.

The v18.0 route computes the \emph{antisymmetric} (Lie/Berry) sector of the
self-model quantum geometric tensor (QGT), in contrast with the v17.0 route which
computed the \emph{symmetric} (Fubini--Study/cone-Hessian) sector and returned
\textsc{none}. By the Provost--Vall\'ee split [PV80], these are the imaginary and
real parts of one object $Q$; the v17.0 \textsc{none} therefore binds only the real
part and \emph{does not bind} the Berry verdict computed here.

\begin{definition}[The QGT recipe; convention lock]
\label{def:qgt}
For a Hermitian rank-$1$ projector $P=P(x)$ depending on real parameters $x^\mu$,
\begin{equation}
  Q_{\mu\nu} \;=\; \Tr\!\big(P\,\partial_\mu P\,\partial_\nu P\big),
  \qquad
  g_{\mu\nu} \;=\; \Re Q_{\mu\nu}\ \text{(Fubini--Study)},
  \qquad
  \FB{}_{\mu\nu} \;=\; -2\,\Im Q_{\mu\nu}\ \text{(Berry curvature)} .
  \label{eq:qgt}
\end{equation}
\end{definition}

Two design points are load-bearing and must not be gotten wrong.

\begin{remark}[Design point 1: the associative complex product, not the Jordan product]
\label{rem:assoc}
The projector $P$ is realized as the complex matrix $P=\texttt{slice\_to\_complex}(E(x))$,
in which the $\Cu$ entries $a_0+a_7 e_7$ are mapped to $a_0+a_7 i$ ($e_7\mapsto i$,
$i=\texttt{sympy.I}$ \emph{symbolic}). The products in \eqref{eq:qgt} are the
\emph{associative} complex matrix products, \emph{not} the Jordan product
$\circ=\tfrac12(AB+BA)$. This is forced: $\Cu\cong\mathbb{C}$ is the unique associative
completion that $u=e_7$ provides, so the associative operator product is canonical,
not chosen. Using $\circ$ inside $P\,\dd P\,\dd P$ would compute a different, wrong
object (forbidden proxy \texttt{fp-octonion-algebra}). Consequently $\FB$ is the
\emph{literal} SymPy imaginary part of $Q$.
\end{remark}

\begin{remark}[Design point 2: the real-part anchor is soft]
\label{rem:soft}
The comparison $\Re(Q)$ vs.\ the v17.0 cone-Hessian (CALC-03,
\S\ref{sec:calc03}) is a \emph{soft, informative diagnostic only}. A discrepancy is
\emph{expected} --- the Fubini--Study pullback is not the cone-Hessian restriction
--- and is \emph{never} a kill or halt (forbidden proxy
\texttt{fp-reuse-cone-hessian}). The v17.0 \textsc{none} binds the real part; this
phase mines the imaginary part.
\end{remark}

\section{The $\mathbb{CP}^1$ convention pin}
\label{sec:cp1}

Before constructing any $\HO$ object we calibrate the sign and the factor of $2$ in
\eqref{eq:qgt} against the canonical magnetic monopole on
$\mathbb{CP}^1=S^2$. With
\begin{equation}
  P(\theta,\phi)=\frac12
  \begin{pmatrix} 1+\cos\theta & \sin\theta\,e^{-i\phi}\\[2pt]
                  \sin\theta\,e^{i\phi} & 1-\cos\theta \end{pmatrix},
  \qquad P^2=P,\ \ P=P^\dagger,\ \ \Tr P=1,
\end{equation}
the recipe \eqref{eq:qgt} gives, exactly,
\begin{equation}
  Q_{\theta\theta}=\tfrac14,\qquad
  Q_{\phi\phi}=\tfrac14\sin^2\theta,\qquad
  Q_{\theta\phi}=\tfrac{i}{4}\sin\theta,
  \label{eq:cp1Q}
\end{equation}
so that
\begin{equation}
  g_{\theta\theta}=\Re Q_{\theta\theta}=\tfrac14,
  \qquad
  \FB{}_{\theta\phi}=-2\,\Im Q_{\theta\phi}=-\tfrac12\sin\theta,
  \qquad
  \int_{S^2}\FB{}_{\theta\phi}\,\dd\theta\,\dd\phi=-2\pi .
  \label{eq:cp1FB}
\end{equation}
The first-Chern number is $c_1=\tfrac{1}{2\pi}\int_{S^2}\FB=-1$, the standard
$\mathbb{CP}^1$ monopole. Equations \eqref{eq:cp1Q}--\eqref{eq:cp1FB} pin the sign and
the factor of $2$ of the entire recipe and confirm that the SymPy $\Re/\Im$ split
behaves as required. (Verified exactly; \texttt{cartan\_phaseA5\_berry.py},
\texttt{cp1\_pin}.)

\section{The rank-1 $\Cu$ idempotent family and the $e_7\mapsto i$ bridge}
\label{sec:family}

On the $\Cu$ chart take
\begin{equation}
  v=(1,z_1,z_2)^{\!\top},\qquad z_1=a+ib,\ \ z_2=c+id,
  \qquad
  P_{\mathbb{C}}=\frac{vv^\dagger}{v^\dagger v},
  \label{eq:family}
\end{equation}
with $a,b,c,d$ real. Exactly over $\mathbb{Q}(i)$ one verifies $P_{\mathbb{C}}^2=P_{\mathbb{C}}$,
$P_{\mathbb{C}}=P_{\mathbb{C}}^\dagger$, $\Tr P_{\mathbb{C}}=1$. This $P_{\mathbb{C}}$ \emph{is}
$\texttt{slice\_to\_complex}$ of the rank-$1$ $\Cu$ idempotent: at any rational test
point the hand-built $\HCu$ element (with off-diagonal $\Cu$ entries reconstructed
from the lower-triangular complex entries of $P_{\mathbb{C}}$) maps under
$\texttt{slice\_to\_complex}$ to exactly $P_{\mathbb{C}}$, and no octonion component
$e_1,\dots,e_6$ survives (the $\Cu$ projection is faithful; if any did, the
projection would be mis-implemented and the verdict untrustworthy). The four
differentiation directions $(a,b,c,d)$ are the $\Cu^2$ survivors $\{11,18,19,26\}$ of
Phase~75. (Verified; \texttt{task1\_idempotent\_and\_bridge}.)

The $4\times4$ QGT $Q(a,b,c,d)=\big[\Tr(P_{\mathbb{C}}\,\partial_\mu P_{\mathbb{C}}\,
\partial_\nu P_{\mathbb{C}})\big]_{\mu,\nu\in(a,b,c,d)}$ assembles over $\mathbb{Q}(i)$ with
$\Re Q$ symmetric and $-2\Im Q$ antisymmetric by construction. At the base point
$a=b=c=d=0$ ($v=(1,0,0)^\top$, $P=\mathrm{diag}(1,0,0)$),
\begin{equation}
  Q(0)=\begin{pmatrix}
    1 & i & 0 & 0\\ -i & 1 & 0 & 0\\ 0 & 0 & 1 & i\\ 0 & 0 & -i & 1
  \end{pmatrix},
  \qquad
  g(0)=\mathbb{1}_4,
  \qquad
  \FB(0)=\begin{pmatrix}
    0 & -2 & 0 & 0\\ 2 & 0 & 0 & 0\\ 0 & 0 & 0 & -2\\ 0 & 0 & 2 & 0
  \end{pmatrix}.
  \label{eq:Qbase}
\end{equation}

\section{VALD-03: $\FB$ is well-defined, generically nonzero, born from breaking}
\label{sec:vald03}

\subsection{Born from the $\Cu$ breaking (group theory)}

\begin{proposition}[Round Berry vanishes; the source is the $\Cu$ breaking]
\label{prop:born}
$\OP=F_4/\Spin(9)$ is isotropy-irreducible: the isotropy representation is the unique
irreducible $16$-dimensional real $\Spin(9)$ spinor representation [Baez, Sec.~3.4].
An isotropy-irreducible homogeneous space admits a unique invariant symmetric
bilinear form up to scale (the Fubini--Study metric) and \emph{no} invariant
$2$-form, because an invariant $2$-form is a $\Spin(9)$-invariant element of
$\Lambda^2 V^*$ for $V$ the (real-type) irreducible spinor, of which there is none
[Berger]. Hence the round ($F_4$-invariant) Berry curvature vanishes identically.
Therefore any nonzero $\FB$ is \emph{born from the $\Cu$ breaking}: the
$\Cu$-reduced isotropy admits the K\"ahler $2$-form, and that is the source of
$\FB$.
\end{proposition}

This is the load-bearing structural fact: the Berry $2$-form cannot be an artifact of
the round geometry --- if it is nonzero it is a genuine consequence of selecting the
complex line $\Cu\subset\mathbb{O}$.

\subsection{Generically nonzero (exact computation)}

Computing $\FB=-2\Im Q$ from \eqref{eq:family} symbolically over $\mathbb{Q}(i)$:

\begin{lemma}[$\FB$ is well-defined and generically nonzero]
\label{lem:nonzero}
$\FB(a,b,c,d)$ is not the identically-zero $4\times4$ matrix. At the base point it is
the block-diagonal $\mathbb{CP}^2$ K\"ahler form $\FB(0)$ of \eqref{eq:Qbase}
($\FB{}_{ab}=\FB{}_{cd}=-2$, off-blocks $0$), and it remains nonzero and antisymmetric
at a generic rational point, e.g.\ $(a,b,c,d)=(\tfrac12,-\tfrac13,2,\tfrac15)$. A
representative symbolic entry,
\begin{equation}
  \FB{}_{ab}(a,b,c,d)=\frac{-2(c^2+d^2+1)}{\big(1+a^2+b^2+c^2+d^2\big)^2},
  \label{eq:FBsym}
\end{equation}
is a nonzero rational function on the slice. (Verified exactly; \texttt{vald03}.)
\end{lemma}

\begin{remark}[Degeneracy guard; negative-result-is-success]
\label{rem:degen}
Had $\FB$ been identically zero on the slice it would be the \emph{degenerate}
(worse-than-soft-kill) outcome --- no antisymmetric content at all --- and would be
reported \emph{flat}, never relabelled ``effectively present'' (forbidden proxy
\texttt{fp-relabel-vacuum}). It is not: $\FB{}_{ab}=\FB{}_{cd}=-2\neq0$.
\end{remark}

\section{CALC-03: the soft real-part anchor (informative, not a kill)}
\label{sec:calc03}

The real part $g=\Re Q=\tfrac12\Tr(\dd P\,\dd P)$ is a Fubini--Study metric. At the
base point $g(0)=\mathbb{1}_4$ (the round $\mathbb{CP}^2$-over-$\Cu$ metric) and it is
positive-definite there and at a generic rational point (eigenvalues $>0$ exactly over
$\mathbb{Q}$). This positive-definiteness is the \emph{only} hard check in CALC-03: a
non-positive-definite $\Re Q$ would mean the QGT construction or the family is wrong,
and would force a stop before trusting $\Im Q$. It passes.

The non-Einstein \emph{character} of $g$ is then compared --- informatively --- to the
v17.0 cone-Hessian. The warm engine reproduces the cone-Hessian
$\mathrm{diag}(9,9,18,18)$ (eigenvalue pattern $9^{(2)},18^{(2)}$) and the round $H^3$
curvature $K=-1$ ($R=-6$; the Phase-71 cone-Hessian slice has $K=-\tfrac12$). The
round Fubini--Study $\Re Q$ has the maximally-symmetric pattern $1^{(4)}$ (a single
eigenvalue of multiplicity $4$: constant holomorphic sectional curvature), which
differs from the cone-Hessian's $2+2$ split.

\begin{remark}[\texttt{fp-reuse-cone-hessian} guard]
\label{rem:fpcone}
This discrepancy is \emph{expected}: the Fubini--Study pullback of the rank-$1$
projector family is not the cone-metric Hessian $\mathrm{Hess}(-\log\det)$ restricted
to the slice [FK94]; they are different objects. The comparison is reported as a
diagnostic and is \emph{never} a kill or halt on this anchor alone. The v17.0
\textsc{none} verdict is a statement about the real (cone-Hessian) sector and is
\emph{not} load-bearing for the imaginary (Berry) sector decided here; treating it as
load-bearing would re-bind to the v17.0 \textsc{none} by sector confusion. The Re
(symmetric) block is type-separated from the Im (antisymmetric) verdict.
\end{remark}

\section{CALC-04: the $M=0$ vacuum class of $\FB$}
\label{sec:calc04}

We classify the vacuum $\FB$ exactly. The $\Cu$ complex structure $J$ on the
$4$-dimensional tangent space pairs $a\leftrightarrow b$ and $c\leftrightarrow d$
(from $z_1=a+ib$, $z_2=c+id$), with $J^2=-\mathbb{1}$. The K\"ahler form is
$\omega_K{}_{\mu\nu}=g(J e_\mu,e_\nu)$; at $M=0$, $g(0)=\mathbb{1}_4$ so $\omega_K$ is
the standard symplectic form pairing $a\leftrightarrow b$, $c\leftrightarrow d$.

The natural trace/traceless/Weyl analog for a $2$-form is its Lefschetz split into a
\emph{primitive-scalar} (``trace'') part proportional to $\omega_K$ and a
\emph{primitive} (``traceless'') remainder:
\begin{equation}
  \FB \;=\; s\,\omega_K \;+\; \big(\FB - s\,\omega_K\big),
  \qquad
  s=\frac{\langle \FB,\omega_K\rangle}{\langle \omega_K,\omega_K\rangle},
  \qquad
  \langle A,B\rangle=\tfrac12\textstyle\sum_{\mu\nu}A_{\mu\nu}B_{\mu\nu}.
\end{equation}

\begin{proposition}[The $M=0$ vacuum is pure-$\Lambda$/K\"ahler]
\label{prop:vacuum}
At $M=0$, $s=-2$ and the primitive remainder $\FB-s\,\omega_K$ vanishes exactly.
Thus the vacuum Berry curvature is pure-K\"ahler,
\begin{equation}
  \FB^{\mathrm{vac}} \;=\; -2\,\omega_K,
\end{equation}
the maximally-symmetric $\mathbb{CP}^2$ K\"ahler $2$-form. It is not flat ($\FB\neq0$)
and carries no ``other'' (primitive) content. (Verified exactly; \texttt{calc04}.)
\end{proposition}

\begin{remark}[\texttt{fp-relabel-vacuum} guard; Pitfall 4]
\label{rem:fprelabel}
The constant $-2\,\omega_K$ is the $\Lambda$/K\"ahler \emph{vacuum} $2$-form, not
matter curvature: matter must vanish as $M\to0$ (that is the burden of plan 76-02).
The cosmological constant is \emph{not} reintroduced negative, and the falsified v17.0
$\mathbb{R}\times H^3$ vacuum (CONVENTIONS \S6) is \emph{not} reintroduced.
\end{remark}

\paragraph{Hand-off to plan 76-02 (pre-registered).} The vacuum Berry curvature
$\FB^{\mathrm{vac}}=-2\,\omega_K$ (base value in \eqref{eq:Qbase}) and the symbolic
slice form $\FB(a,b,c,d)$ (\eqref{eq:FBsym} and companions, computed by
$\texttt{berry\_F}(P_{\mathbb{C}},(a,b,c,d))$) are emitted so that plan 76-02 may
\emph{subtract} the vacuum and isolate the matter contribution
$\FB^{\mathrm{matter}}:=\FB(M)-\FB^{\mathrm{vac}}$, which must vanish as $M\to0$.

\section{Forbidden proxies explicitly rejected (PART 1)}
\label{sec:fp}

\begin{itemize}
\item \texttt{fp-reuse-cone-hessian} --- rejected (Remark~\ref{rem:fpcone}): the
  cone-Hessian / $\Re(Q)$ is the real-part soft consistency check only and is never
  load-bearing for the imaginary-part verdict; no kill or halt is issued because
  $\Re(Q)$ differs from the cone-Hessian character.
\item \texttt{fp-relabel-vacuum} --- rejected (Remarks~\ref{rem:degen},
  \ref{rem:fprelabel}): an identically-zero $\FB$ would be reported flat (it is not);
  the constant $-2\,\omega_K$ is named the $\Lambda$/K\"ahler vacuum, not matter; no
  $\Lambda<0$ / $\mathbb{R}\times H^3$.
\item \texttt{fp-float-decisive} --- rejected: the imaginary unit is
  \texttt{sympy.I} (symbolic, $e_7\mapsto i$); every decisive quantity ($\FB$
  non-vanishing, $\Re Q$ positive-definiteness, the vacuum class, the $\mathbb{CP}^1$
  pin) is rational/symbolic over $\mathbb{Q}(i)$ via SymPy; no numpy ranks or float
  eigenvalues are used.
\item \texttt{fp-octonion-algebra} --- rejected (Remark~\ref{rem:assoc}):
  \texttt{octonion\_algebra.py} is absent on the decisive path (runtime source guard);
  det/trace/Jordan primitives come from \texttt{ring\_lemma\_verification.py}; the QGT
  uses the associative complex matrix product on $\texttt{slice\_to\_complex}(P)$,
  never the Jordan product.
\end{itemize}

% ============================================================================
% ============================================================================
%   PART 2  (plan 76-02, the DECISIVE / MATTER half)
% ============================================================================
% ============================================================================
% Backed exact-over-Q(i) by code/cartan_phaseA5_berry.py PART 2 (ALL_PASS, exit 0):
%   matter-on E(x;M) over C_u (rank-1 + C_u-faithful + M->0); the matter F_B M-series;
%   the F_B^F_B Lorentz-block epsilon-contraction; the SO(3,1) frame-rotation test; the
%   independently-frozen T[M] (psi=2Re((x2 x1)x3), kappa first, off-switch, no F_B);
%   the VALD-04 Einstein-vs-EM verdict on gauge-invariant scalars.
% Additional references for PART 2:
%   [Wise]  D.K. Wise, "MacDowell-Mansouri gravity and Cartan geometry", gr-qc/0611154
%             (2010). The epsilon-contraction F^F -> EH + Lambda: the F_B^F_B Lorentz-
%             block contraction is the A.5(b) DIAGNOSTIC (the contraction STRUCTURE only,
%             NOT the forced-vs-posited Phase-C audit).
%   [MM77]  S.W. MacDowell, F. Mansouri, Phys. Rev. Lett. 38 (1977) 739. The original
%             unification of the Lorentz connection and coframe into one curvature whose
%             F^F gives EH+Lambda -- the structural template for the F_B^F_B contraction.

\bigskip
\begin{center}\rule{0.6\textwidth}{0.4pt}\end{center}
\bigskip

\part*{PART 2 (plan 76-02) --- the decisive / matter half}
\addcontentsline{toc}{section}{PART 2 (plan 76-02) --- the decisive verdict}

\renewcommand{\thesection}{B.\arabic{section}}
\setcounter{section}{0}
\renewcommand{\theproposition}{B.\arabic{proposition}}
\renewcommand{\thelemma}{B.\arabic{lemma}}
\renewcommand{\theremark}{B.\arabic{remark}}

PART~1 fixed the vacuum baseline exactly over $\mathbb{Q}(i)$: the QGT recipe is
calibrated, $\FB=-2\,\Im Q$ is well-defined and generically nonzero (VALD-03), $\Re Q$
is a positive-definite Fubini--Study metric (CALC-03, soft), and the $M=0$ vacuum is
pure-$\Lambda$/K\"ahler $\FB^{\mathrm{vac}}=-2\,\omega_K$ (CALC-04). PART~2 turns matter
$M\in V_{1/2}$ \emph{on} and decides the genuinely-open VALD-04 clause on
gauge-invariant scalars: is the matter-sourced Berry curvature transverse /
Einstein-shaped (M-power, tensor structure, \emph{and} support matchable to an
independently-frozen $V_{1/2}$ stress-energy $T[M]$ $\Rightarrow$ \textsc{survives}),
or EM-shaped (a traceless $\sim F^2$ field strength, structurally disjoint from $T[M]$
$\Rightarrow$ \textsc{soft kill}, reproducing the v17.0 same-wall mismatch on the Lie
sector)?

\section{The matter-on idempotent $E(x;M)$ over $\Cu$ (the biggest modeling gap)}
\label{sec:matter-idempotent}

Matter in $V_{1/2}(E_{11})$ (engine indices $\{11..26\}$) is the octonionic content of
$x_2$ (entry $(0,2)$, $\{11..18\}$) and $x_3$ (entry $(1,0)$, $\{19..26\}$). Under the
$\Cu$ bottleneck $\pi_u$ (the $e_7$-projection) the matter splits cleanly into two
species, and \emph{this split is the resolution of the $\Cu$-faithfulness question}:

\begin{proposition}[$\Cu$-faithfulness of the matter family]
\label{prop:faithful}
The $\Cu$-survivor matter $\{11,18,19,26\}=\{\Re x_2,\langle x_2,e_7\rangle,\Re x_3,
\langle x_3,e_7\rangle\}$ survives $\pi_u$ and is a genuine complex off-diagonal
perturbation of the $\Cu$ idempotent; the $\Cu$-transverse matter $\{12..17,20..25\}$
(the $e_1..e_6$ components of $x_2,x_3$) is annihilated by $\pi_u$
($\mathrm{proj}_u(e_k)=0$, $k=1..6$). Consequently the matter-on idempotent
$E(x;M)=v v^\dagger/(v^\dagger v)$ over $\Cu$, with the matter shifting the Veronese
centre,
\begin{equation}
  z_1 = a + i b + \mu_1 t\ \ (x_3\text{ slot}),\qquad
  z_2 = c + i d + \mu_2 t\ \ (x_2\text{ slot}),
  \label{eq:matter-chart}
\end{equation}
is a clean rank-$1$ $\Cu$ projector ($E^2-E=0$, $\Tr E=1$, $E=E^\dagger$) to \emph{all}
orders in the matter amplitude $t$, and the transverse matter is a projection (not a
rank-changing leak). Verified exactly over $\mathbb{Q}(i)$.
\end{proposition}

The transverse-matter projection is not a pathology: it is the statement that the
canonical $\Cu$ Berry curvature only \emph{sees} the $\Cu$-faithful part of the matter
(the part Phase~75 already isolated as the $4$ coframe directions). No Approach-2
octonionic cross-check is triggered, because the family is a clean $\Cu$ projector.

\begin{lemma}[The $M\to0$ limit and the leading $M$-power]
\label{lem:Mlimit}
The matter Berry curvature $\FB^{\mathrm{matter}}:=\FB(M)-\FB^{\mathrm{vac}}$ vanishes
exactly at $t=0$ (the PART-1 vacuum $\FB^{\mathrm{vac}}=-2\,\omega_K$ is recovered), and
its leading behaviour is \emph{second order} in the matter amplitude,
\begin{equation}
  \FB^{\mathrm{matter}}_{\mu\nu} = O(\|M\|^2),
  \label{eq:Mpower}
\end{equation}
confirmed by the exact ratio test $\FB^{\mathrm{matter}}(t)/\FB^{\mathrm{matter}}(t/2)
\to 4 = 2^2$ (and $\FB^{\mathrm{matter}}/t\to0$, $\FB^{\mathrm{matter}}/t^2\to$const) as
$t\to0$, at two independent rational matter directions. This is \emph{exactly} the order
at which the v17.0 metric perturbation $h$ entered.
\end{lemma}

\section{The matter $\FB$, the $\FB\wedge\FB$ contraction, and the EM signature}
\label{sec:FBmatter}

The matter Berry curvature is, at $t=\tfrac1{20}$ and the base point,
\begin{equation}
  \FB^{\mathrm{matter}}=\frac{1}{6889}
  \begin{pmatrix}
    0 & 658 & -224 & 32\\
   -658 & 0 & -32 & -224\\
    224 & 32 & 0 & 818\\
   -32 & 224 & -818 & 0
  \end{pmatrix},
  \label{eq:FBmatter}
\end{equation}
manifestly \emph{antisymmetric} --- a genuine $2$-form / field strength.

\subsection{The MacDowell--Mansouri $\FB\wedge\FB$ Lorentz-block contraction}

Following the MM template~[Wise, MM77], the diagnostic scalar is the
$\varepsilon$-contraction of $\FB\wedge\FB$ onto the forced $\SO(3,1)$ Lorentz block of
Phase~75 ($G=\mathrm{diag}(+1,-1,-1,-1)$ on the survivors $\{11,18,19,26\}$):
\begin{equation}
  I_{FF} \;=\; \varepsilon^{\mu\nu\rho\sigma}\,\FB_{\mu\nu}\,\FB_{\rho\sigma}.
  \label{eq:IFF}
\end{equation}
This is the Pontryagin / instanton density of the $2$-form $\FB$, a gauge-invariant
scalar. Exactly over $\mathbb{Q}(i)$: $I_{FF}^{\mathrm{matter}}=46944/571787$,
$I_{FF}^{\mathrm{full}}=16384000/571787$, $I_{FF}^{\mathrm{vac}}=32$ --- all nonzero. In
the MM construction the $\varepsilon$-contraction of the \emph{curvature} of the Cartan
connection produces the Einstein--Hilbert plus cosmological terms; for a pure
\emph{field strength} $\FB$ it is the topological $\sim F^2$ density.

\subsection{The decisive tensor-structure discriminant: the Maxwell stress is traceless}

The field-strength (Maxwell) stress of $\FB$ with respect to the forced Lorentz metric,
\begin{equation}
  T[\FB]_{\mu\nu} \;=\; \FB_{\mu a}\,G^{ab}\,\FB_{\nu b}
    \;-\; \tfrac14\,G_{\mu\nu}\,F^2,\qquad
  F^2 = G^{ac}G^{bd}\FB_{ab}\FB_{cd},
  \label{eq:maxwell}
\end{equation}
satisfies, \emph{exactly over} $\mathbb{Q}(i)$ and in $n=4$,
\begin{equation}
  \boxed{\;G^{\mu\nu}\,T[\FB]_{\mu\nu} \;=\; 0\;}
  \qquad(F^2 = 472320/47458321\neq0).
  \label{eq:traceless}
\end{equation}
This vanishing trace is the \emph{conformal invariance of any antisymmetric $2$-form in
exactly four dimensions} (an $n=3$ control gives trace $=2\neq0$, so~\eqref{eq:traceless}
is a genuine $4d$ fact, not a coincidental zero): $\FB$ is a conformal $\sim F^2$ field
strength, \emph{not} a non-traceless Einstein stress. This is the decisive
tensor-structure discriminant.

\section{The $\SO(3,1)$ frame-rotation gauge test (fp-nonabelian-gauge)}
\label{sec:frame}

Although the rank-$1$ idempotent makes $\FB$ abelian, the Lorentz-block contraction and
every frame statement are tested under a generic $\SO(3,1)$ frame rotation $L$ (a
rational boost $\cosh=5/4$ composed with a rational rotation $\cos=3/5$;
$L^\top G L = G$ exactly over $\mathbb{Q}$). The Pontryagin scalar~\eqref{eq:IFF} is
\emph{unchanged} ($I_{FF}=46944/571787$ before and after), and the Maxwell
trace~\eqref{eq:traceless} stays exactly zero. The verdict scalars are gauge-invariant,
not frame artifacts.

\section{The independently-frozen stress-energy $T[M]$ (kappa first, no $\FB$)}
\label{sec:TM}

Following the Phase-73 DERV-03 discipline, the matter $\FB$ is compared against an
\emph{independently-frozen} $V_{1/2}$ stress-energy built from the
$V_0\leftrightarrow V_{1/2}$ cross-term content \emph{without any reference to} $\FB$
(non-circular; $\kappa$ frozen first; AST-guarded so no $\mathrm{Ric}/R/G$ symbol enters
its construction). With $\psi(x;M)=2\,\Re\!\big((x_2 x_1)x_3\big)$ the certified
$\det_3$ SSOT cross-term scalar (built with $\mathtt{oct\_mul}$ in the SSOT order
$(x_2 x_1)x_3$, never $\mathtt{octonion\_algebra.py}$),
\begin{equation}
  T[M]_{\mu\nu}=T[\psi]_{\mu\nu}
    =\partial_\mu\psi\,\partial_\nu\psi
     -\tfrac12\,\eta^{\mathrm{bg}}_{\mu\nu}\,(\partial\psi)^2,
  \qquad
  \mathrm{tr}_\eta T[M]=680\neq0.
  \label{eq:TM}
\end{equation}
$T[M]$ is \emph{non-traceless} --- the Einstein-source tensor signature, structurally
distinct from the traceless Maxwell stress~\eqref{eq:traceless} of $\FB$. The
\emph{off-switch} certifies $T[M]$ is genuinely cross-term-sourced: the matter-sourced
spacetime curvature changes decisively when the $V_0\leftrightarrow V_{1/2}$ triple is
dropped ($\mathtt{norm\_potential}=\mathtt{inv\_det\_X\_block}$),
$R[g]_{\mathrm{ON}}\approx5.13$ vs $R[g]_{\mathrm{OFF}}\approx1144.87$, and its
$n=4$ Ricci decomposition is genuinely non-trivial ($R\neq0$, $S\neq0$, Weyl $\neq0$).
This is the v17.0 \emph{symmetric}-sector object: it \emph{builds} $T[M]$ and is
\textbf{never} the imaginary-part verdict (fp-reuse-cone-hessian).

\section{VALD-04: the three matches and the verdict}
\label{sec:verdict}

The verdict is the three-way match of the matter $\FB$-source against the
independently-frozen $T[M]$, on gauge-invariant scalars:
\begin{enumerate}
\item \textbf{$M$-power:} $\FB^{\mathrm{matter}}\sim O(\|M\|^2)$ and
  $\mathrm{tr}_\eta T[M]\sim O(\|M\|^2)$ ($\psi\sim$ matter, $(\partial\psi)^2\sim$
  matter$^2$). The powers \emph{coincide}.
\item \textbf{Tensor structure (decisive mismatch):} $\FB$ is an antisymmetric $2$-form
  whose Maxwell stress is \emph{exactly traceless} ($\sim F^2$, conformal) ---
  \textbf{EM-shaped}; $T[M]$ is symmetric and \emph{non-traceless} ---
  Einstein-source. These are different tensor species.
\item \textbf{Support:} a traceless conformal $\sim F^2$ stress cannot equal a
  non-traceless gradient stress; the supports / structures are \emph{disjoint}.
\end{enumerate}

\begin{proposition}[VALD-04 verdict --- \textsc{soft kill}]
\label{prop:verdict}
The matter Berry curvature $\FB^{\mathrm{matter}}$ is \textbf{EM-shaped}: an
antisymmetric $2$-form whose $\FB\wedge\FB$ Lorentz-block contraction is the Pontryagin
$\sim F^2$ density and whose Maxwell stress is exactly traceless (conformal,
\eqref{eq:traceless}), \emph{structurally distinct} from the non-traceless, symmetric,
independently-frozen $T[M]$~\eqref{eq:TM}. Although the imaginary (Lie / antisymmetric)
sector is a \emph{different tensor} than the v17.0 symmetric cone-Hessian, it
\textbf{inherits the symmetric-sector same-wall mismatch}: the matter $\FB$ does not
match $T[M]$ in tensor structure or support. Decided on gauge-invariant scalars,
frame-rotation invariant, exact over $\mathbb{Q}(i)$.
\end{proposition}

\begin{center}\fbox{\parbox{0.92\textwidth}{\centering
\textbf{Phase A.5: \textsc{soft kill}.} The matter-sourced canonical Berry curvature is
EM-shaped (an antisymmetric field strength, traceless $\sim F^2$ Maxwell stress,
structurally disjoint from the independently-frozen $T[M]$), \emph{not} Einstein-shaped.
The Lie sector inherits the v17.0 symmetric-sector same-wall mismatch. \textbf{Recommend
STOP before Phase~B.} Phase~A.5 is \emph{conjunctive} with Phase~75 (\textsc{survives}):
since A.5 is a \textsc{soft kill}, the conjunction $\big(\text{Ph75}=\textsc{survives}\big)
\wedge\big(\text{Ph76}=\textsc{soft kill}\big)$ does \textbf{not} greenlight Phase~77;
v18.0 closes as a \emph{publishable negative}. Reported flat at true strength
(negative-result-is-success).
}}\end{center}

\section{Forbidden proxies explicitly rejected (PART 2)}
\label{sec:fp2}

\begin{itemize}
\item \texttt{fp-relabel} --- rejected: the verdict is the three-way ($M$-power $+$
  tensor structure $+$ support) match against an \emph{independently-frozen} $T[M]$
  (built without $\FB$, $\kappa$ first), \emph{never} ``a $2$-form appears''. A nonzero
  antisymmetric $\FB$ is a \emph{field strength} (EM, source side), not an Einstein
  structure --- exactly the v17.0 Ph73 lesson.
\item \texttt{fp-reuse-cone-hessian} --- rejected: the cone-Hessian / $\Re(Q)$ and the
  symmetric spacetime curvature $R[g]$ (with $S\neq0$, Weyl $\neq0$) are used \emph{only}
  to build the independently-frozen $T[M]$; they are never load-bearing for the
  imaginary-part verdict. The v17.0 \textsc{none} binds only $\Re$.
\item \texttt{fp-relabel-softkill} --- rejected: the EM-shaped result is reported
  \emph{flat} as a \textsc{soft kill}, never ``approximately Einstein''; an
  Einstein-shaped \textsc{survives} would not have been deflated (claims at true strength
  in both directions).
\item \texttt{fp-nonabelian-gauge} --- rejected: the verdict rests on gauge-invariant
  scalars (the $\FB\wedge\FB$ Pontryagin scalar, the Maxwell trace) tested and confirmed
  invariant under a generic $\SO(3,1)$ frame rotation.
\item \texttt{fp-float-decisive} --- rejected: every decisive number is exact over
  $\mathbb{Q}(i)$ via SymPy ($\Im$ literal; the Maxwell trace is \emph{exactly} $0$, the
  $T[M]$ trace \emph{exactly} $680$, the Pontryagin scalar an exact rational); no float.
\item \texttt{fp-octonion-algebra} --- rejected: \texttt{octonion\_algebra.py} is absent
  on the decisive path (runtime source guard); $\psi$ uses $\mathtt{oct\_mul}$ in the
  $\det_3$ SSOT order; the QGT uses the associative complex matrix product on
  $\mathtt{slice\_to\_complex}(P)$, never the Jordan product.
\end{itemize}

\end{document}
